\documentclass[12pt]{article}

\usepackage{hyperref}
\usepackage{natbib}
\bibpunct[: ]{(}{)}{;}{a}{}{,}
\usepackage[margin=1in]{geometry}
\usepackage{setspace} 
\usepackage{graphicx}
\usepackage{lscape}
\usepackage{subfigure}

\hypersetup {
pdfauthor={Omar Lizardo},
pdfpagemode = UseNone,
pdfstartview=FitH,
citebordercolor = 1 0 0,
}

\begin{document}

\title{Patterns of cultural choice and the instrumental use of social ties}

\author{Omar Lizardo \\ \textit{University of Notre Dame}}


\date{\today}
\maketitle

\begin{abstract}
Recent research in the sociology of culture has argued and provided empirical evidence for the general proposition that persons who have diverse tastes should be more likely to also have larger personal networks, and should also be more likely to have social ties that span relatively larger distances in social space. A different strand of network theory focuses not only on relative differences in the size and ``reach'' of personal networks but also on the relative opportunities and resources that are made available to individuals by virtue of having those ties. In this paper, I draw a theoretical link between cultural capital, in the form of a command of a variety of cultural resources and the instrumental use of social ties. I build on the general connection between culture consumption and sociability which has been the focus of recent research in the sociology of culture. I propose a model in which the possession of cultural resources leads to a higher likelihood of activation of networks ties. The results show that net of other individual predictors, as the cultural variety of the person's repertoire of tastes increases, there is 1) a higher likelihood of activating social ties to obtain information about jobs, and 2) a higher likelihood that the tie activation is through a weak, rather than a strong tie which 3) leads them to be less likely to experience long unemployment
spells. 
\end{abstract}


\doublespacing
\newpage 
\section{Introduction}
Recent research in the sociology of culture has argued and provided empirical evidence for the general proposition that culture consumption should be connected to networks of larger size, and that cultural variety, especially of cultural aptitudes that are widely accessible and are spread fairly evenly across the social structure, is associated with the individual possession of social ties that span relatively larger distances in social space. For instance, \citet[247]{erickson96} notes that ``[t]he more diverse the set of{\ldots}contacts one has, the more variety of culture one will encounter and hence maintain or learn{\ldots}Thus network variety is one important source of cultural variety.'' DiMaggio proposes that ``[t]he reason for this finding, so at odds with conventional notions about the isomorphism of taste and class{\ldots}, is that wide-ranging networks require broad repertoires of taste'' \citeyearpar[444]{dimaggio87}. \citet{lizardo06a} highlights the additional process of how tastes are \emph{converted} into a wider range of network ties. Thus, in the same way that network variety leads to cultural variety, having a wide-ranging repertoire of tastes is necessary to \emph{sustain} large networks over time, especially in the face of the large-scale volatility of personal relationships among high status persons, who are increasingly subject to pressures related to geographical mobility and other lifestyle choices that serve to disrupt current network relations.

This strand of research however, has so far ignored the fact that network theory focuses not only on relative differences in the size and density of individual networks---indexed by positional measures such as the various centrality indices \citep{freeman77, borgatti05}---but also on the relative \emph{resources} that are made available to individuals by virtue of having those ties \citep{bourdieu86, coleman88, portes98, lin01b}, and thus the differential opportunities to access those resources in times of need. In other words, networks are transformed from a potential store of benefits to actual\emph{ social capital} only in the process of agentic \emph{activation} of social ties \citep{lin99, lin00, lin01a, lin01b,mouw03, yakubovich05, lin_etal81a, lin_etal81b, lin_dumin86}.This activation can occur for the purposes of dealing with practical problems and issues like finding a job \citep{granovetter73} or obtaining social support in the form of instrumental benefits or emotional support \citep{wellman_wortley90}. 

In this paper, I empirically examine the connection between cultural capital, in the form of a command of a \emph{variety of cultural resources } \citep{dimaggio87, erickson96, peterson_rossman08} and network activation. I build on the general link between culture consumption and sociability which has been the focus of recent research in the sociology of culture \citep{bennett_etal99, cardon_granjon05, lizardo06a, jeannotte03} to propose a model in which the possession of cultural resources leads to a higher likelihood of activation of networks ties.

This paper is organized as follows: in the next section, I briefly review the literature on the sources and consequences of network activation. I then follow with a discussion of how we can use the general association between culture consumption and sociability in order to develop some theoretical intuition connecting the possession of variety-driven cultural capital and the likelihood of using social networks in a purposive way. In the next section I provide evidence that indeed culture consumption is positively tied to network activation, by analyzing data on the individual determinants of the prototypical form of network activation that has dominated the literature on social capital: using social networks to learn about job opportunities. The results show that net of other individual predictors, as the cultural variety of the individuals' repertoire of tastes increases, there is 1) a higher likelihood of activating social ties to obtain information about jobs, and 2) a higher likelihood that the tie activation is through a weak, rather than a strong tie. Which leads to 3) a relatively lower likelihood of experiencing long unemployment spells. I close by providing a general assessment of the implications of the results and proposals for future research.

\section{Network ties as ``social capital''}
\subsection{Social capital as network diversity}
There are at least two primary ways in which social ties have been conceived as ``valuable'' to individual in contemporary network theory and thus as a source of ``social capital.'' The first has to do with the benefits that can be accrued from having a wide \emph{range} and \emph{diversity} of contacts \citep{burt82, burt83, burt92, lin_dumin86}. The conception of contact diversity as ``social capital'' relies on the well-established observation, first formulated in Granovetter's \citeyearpar{granovetter73} classic paper, that contact homogeneity in sociodemographic traits is associated with information \emph{redundancy}, while contact heterogeneity is associated with access to \emph{novel} sources of information. Whenever instrumental pursuits require access to novel sources of information---as when individuals use social contacts to learn about employment opportunities \citep{granovetter95}---those who are connected to diverse information sources will be more likely to learn about opportunities not otherwise accessible through the local network. Because extra-local information is more likely to be accessed through weak ties rather than strong ties, weak ties should provide the individual with an informational advantage not available to those embedded in local cliques. This is the ``strength of weak ties'' argument.

Studies of the association between social position and local network topologysupport the strength of weak ties argument. For instance, a well-established theme in network theory and research concern the long-standing association between socioeconomic status and the social diversity of contacts to whom the individual is connected. It is a well-established fact that higher status individuals are more likely to have social connections that span a wider distance in social space than low status individuals \citep{lin01a, lin01b}. An extensive variety of empirical evidence (mostly obtained through empirical examination of the ``core discussion network'') supports this contention. We know that the higher educated tend to have larger and less tightly knit (more likely to consist of persons not very well-acquainted with one-another) personal networks than the less educated \citep[129]{marsden87}.

Persons of higher education also tend to enjoy more diverse networks along such dimensions as gender and age \citep{fischer82}. Diversity in the local networks of high-status persons even extend to the \emph{kin network},with high status individuals displaying more heterogeneous kin networks (as measured by the standard deviation of educational attainment and occupational standing) than low status individuals \citep{goldstein_warren00}. In addition, a persons' occupational status (measured by scalar prestige indices) is positively associated both with the \emph{number} of social positions (defined as a detailed occupational classification) that the person is connected to (by having a friend in that job) as well as the \emph{range} of positions (as defined by the standard deviation of occupational prestige across the statuses) contained in their personal network (\citealp[90]{lin01b}; \citealp{lin_dumin86}). That is those of higher status are not only more likely to be connected to other persons of higher status, but are more likely to have social ties with persons of dissimilar status than are low status individuals.

\subsection{Social capital as access to resources}
The second way in which social connections have been conceived as ``social capital'' in network theory has to do with a factor that should be correlated with network diversity: the probability of having access to contacts that control or have access to desirable resources, such as influence, or who occupy high status positions.Thus, a big part of a person's status attainment process occurs through a relational pathway which partly mediated the resources obtained from the class of origins (father's or mother's education) and the individual's own attained status in determining current social position \citep{lin01a}.

A lack of local clustering (e.g. personal contacts being closely tied to one another) in the immediate social network is correlated with the \emph{resources}---as measured by the socio-demographic characteristics of alters---embedded in that network \citep{lin01a}. This explains why individual current achieved status is correlated with a lack of clustering, as measured by the range of positions accessible through the network.Thus, the individual's ascribed characteristics (resources from the family of origin) are associated with the pattern of network relations that he or she currently occupies. The volume of resources that corresponds to the individual's status or origin are positively associated with her current network structure and in part accounts for the association of current achieved status and local network structure \citep{lin90, lin01b}. 

In terms of Bourdieu's \citeyearpar{bourdieu86} theory of social capital, this is the process of ``conversion'' of educational and economic capital in the family of origin into current social capital.Social capital is converted into desirable resources by way of the agentic \emph{activation } of existing social ties. Thus, individuals derive advantage from relatively temporally stable qualities of their social network (size, diversity, rage, etc.) when they are able to call upon the resources embedded in that network as a result of situational contingencies (i.e. needing help finding a job). Thus, before being activated social capital remains in ``stasis'' as it where, and it is only by specific instance of network usage that we can get an empirical handle on the concept of social capital.

\section{Cultural capital and social networks}
A key analytical limitation of network theories of social capital is that they take the observed pattern of network connections as largely \emph{exogenous} to the analysis. That is, the fail to ask the question: where do social ties come from? Recent research at the intersection of the sociology of culture and network theory has begun to address this issue. In particular this research argues that insofar as social relationships are mediated through focused episodes of social interaction, and given that interaction rituals in modern societies tend to revolve primarily around the discussion of symbolic and cultural goods \citep{dimaggio87, dimaggio00} then individuals \emph{cultural tastes, competences and life-style practices,} should be a key predictor of both the size and range of personal networks.

The mobilization of cultural aptitudes and knowledge in conversation is clearly implicated in the \emph{capacity }of individual actors to constantly generate new social ties (and thus sustain a large and diverse social network over time). This insight acquires more salience when considered in the context of recent empirical confirmation of the short lived and ephemeral nature of most network connections \citep{bidart_degenne05, bidart_lavenu05, burt00, burt02, suitor_etal97, suitor_keeton97, feld_etal07}. This new strand of research therefore attempts to argue for the origins of the advantages produced by social network from the \emph{conversion} of cultural capital into social capital \citep{bourdieu86}. 

One of the primary ways in which tastes are used to sustain social networks, is by allowing the individual to connect to alters that are both relationally close---when specialist or (narrow) ``niche'' cultural tastes are concerned---and relationally distant---when talking about aptitude for broad-niche, or ``popular'' cultural forms \citep{lizardo06a}. Culture consumption and the ability to decode and appropriate symbolic goods produced are important resources that individuals mobilize in the \emph{construction} and \emph{maintenance} of active social relationships \citep{cardon_granjon05, dimaggio87, fine77}.  For the most part conversational interaction is devoted to discussion and commentary on symbolic goods \citep{fine77, dimaggio87}.This means that ``[m]any social relations are constructed, in different ways, around cultural and recreational activities'' \citep{cardon_granjon05}.This is an insight that harks back to classical process models of friendship formation, in with similarity based on cultural taste and practices \citep{mark98b}playing the role that ``value-homophily'' played in the original formulation \citep{lazarsfeld_merton54}. 

That the origins of new---and reproduction of previously existing---network connections comes from the deployment of cultural knowledge in interaction is also implicit in previous models of the relationship between culture and social structure \citep{carley91, carley95, fine79, mark98a, mark03}. For instance, the Carley-Mark \citep{carley91, mark98a} ``constructural'' model can be interpreted as an elementary schema of how culture can be translated into social connections, and how social structure (the distribution of chances to interact across persons in the system) and cultural structure (the distribution of cultural forms across persons) can be defined in an interdependent manner. The constructural model is based on a simple assumption of similarity (homophily), which postulates that the likelihood of a social tie increases with the cultural similarity between any given dyad (a dynamic process.\footnote{Contemporary network models of taste formation through social network ties \citep{mark98b} can be seen as partially based on the Lazarsfeld-Merton formulation. However, these models raise what was explicitly a secondary dynamic in the original scheme---i.e. mutual ``value-adjustment'' through social interaction of initially different value dispositions \citeyearpar[33]{lazarsfeld_merton54}---to being the only operative, and thus primary dynamic. It is clear however that for Lazarsfeld and Merton ``selective'' processes whereby contacts are sifted based on similarity of attitudinal dispositions similar to those emphasized in the culture-conversion account were considered primary \citeyearpar[22-29]{lazarsfeld_merton54}. In this way the probability that two persons will interact is driven by their cultural similarity. Interaction, in this positive feedback loop, in its turn increases cultural similarity, as individuals exchange their stocks of knowledge with one another.}

We should therefore expect that individuals who command a wide variety of tastes which cut across these well defined boundaries (i.e. those who have been referred to as ``omnivores'' in recent research \citep{peterson92, peterson05, peterson_rossman08}) and who, given the stability assumption, will also be able to ``keep'' that ``tolerant'' taste pattern active over time---will also be more effective in sustaining high rates of social interaction in their---constantly shifting, sparse and loosely bounded---personal communities.Lizardo \citeyearpar{lizardo06a} using data from the 2002 General Social Survey, finds empirical evidence consistent with these theoretical propositions. He finds that there is a monotonic relationship between the average size of the personal network (as measured by a one-shot item which asked the respondent the number of persons that the he or she comes into contact with at least once a year) and the number of cultural activities---associated with both the fine and the popular arts---that the respondent reports having engaged in at least once in the past year. In addition, the results of the analysis show that the number of as the number of cultural activities that the respondent engages in increases so does the expected size of their personal network, even after accounting for simultaneity bias. Thus it is clear that as cultural capital---defined as the capacity to engage in a wide variety of aesthetic pursuits \citep{bryson96, peterson_kern96, ollivier04}---increases, access to larger (and by implication) more diverse social networks also increases.

\section{Empirical implications}
One of the data sources used in this paper (2002 General Social Survey) contains both direct and indirect measures of weak tie mobilization, while the other one (2002 Current Population Survey) only has information on ``strong ties'' versus direct contact by the employer.  Furthemore, while the GSS contains \emph{retrospective} information on those respondents who \emph{have already found a job}, the CPS contains information on the job search strategies of respondents who were currently unemployed when the survey was conducted.  Thus, I explore the empirical implications of the model proposed above both in the case of individuals who have already found a job, and individuals who are currently looking for a job. 

If the consumption of a variety of cultural goods is connected to the size and density of an individual's personal network \citep{dimaggio87, lizardo06a}, and if networks are a resource that is accessed when the individual requires either support, information or the projection of influence, we should expect those individual endowed with a more varied repertoire of tastes, to be able to sustain not only larger personal networks, but also to be relatively more likely to \emph{access} those networks in a time of need. Thus, one empirical implication of the model proposed so far is that cultural capital in the form of cultural variety should be connected with access of personal networks when individuals require information in comparison to non-relational forms of information search.

We should not only expect cultural taste variety to be connected to the use of relational pathways to information and influence.More specifically, we should expect cultural variety to be connected to information access pathways that are \emph{more likely to be realized by the presence of a weak tie rather than a strong tie}. In the literature weak ties have been conceptualized in two different ways: 1) using a \emph{direct} measure of tie strength by counting weak ties as those named as ``acquaintances'' rather than close friends or relatives \citep{granovetter95, mouw03} and 2) using an \emph{indirect} measurement strategy, in which a persons is counted as having used a weak tie if they obtained their job by being contacted by employer directly (without using an employment service or having answered and ad in the paper) \citep{yakubovich05}.

The rationale for the indirect measurement strategy is that weak ties have been shown to serves as conduits of \emph{both} ``information and direct influence'' \citep[415]{yakubovich05}.The most frequent example of direct influence is of course a direct connection between the potential employer and the individual. It follows therefore that any factor that increases the probability of having a network richer in weak ties will increase the chances of any one individual already having a potential employer as an acquaintance \citep{yakubovich05}.  In either case, I hypothesize that if cultural variety and network variety are positively associated \citep{lizardo06a, erickson96}, and if network variety is defined as having a network rich in contacts who are different from ego, and who are therefore more likely to be weak ties, we should expect that cultural variety will increase the chances of individuals having potential employers as acquaintances, and of thus mobilizing those weak ties when necessary. 

\section{Study 1:  2002 General Social Survey (GSS)}
\subsection{Data and Variables}
I use data from the culture and network modules of the 2002 General Social Survey (Davis, Smith and Marsden, 2002). The GSS is administered biannually by the National Opinion Research Center (NORC) to a nationally representative sample of non-institutionalized, English-Speaking American adults. The 2002 wave of the GSS contained a recurring module on participation in the arts (similar to ones fielded in 1993 and 1998), along with a new module related to social networks and social support. The network module contained a question asking the respondent to report the \emph{method} that they used to obtain their last job, which, in conjunction with the self-report measures of cultural participation, makes these data suitable for testing the hypotheses outlined above.

\subsubsection{Cultural variety scale}
Respondents were asked to report whether they had engaged in the following activities during the past year: 1) gone to see a movie in a theater 2) gone to a live performance of popular music like rock, country, or rap 3) gone to a live performance of a non-musical stage play 4) attended a live ballet or dance performance 5) gone to a classical music or opera performance and 6) visited an art museum or gallery and 7) read a novel, poem or play. The variables are coded one if the respondent engaged in that activity in the past year and zero otherwise. As a measure of cultural variety, I create a linear index (Cronbach's $\alpha=0.93$) that equals the total number of cultural activities (ranging from zero to seven) that the respondent reported engaging in the last year.  This is what \citep{warde_etal08} have referred to as ``omnivorousness by volume.'' Previous research validates the criterion validity of this type of measure as an indicator of consumption patterns characteristic of young, upwardly mobile elites \citep[e.g.][]{bryson96, sullivan_katzgerro07, warde_etal08}.

\subsubsection{Network activation indicator}
This is measured through an item that asked the respondents to ``indicate how you first [from whom] found out about work at your present employer'' respondents who were not currently working were asked to answer this question for their last job. A separate category was reserved for people who had never participated in the labor force. I recode the original variable into a polytomous item containing seven categories: 

\begin{enumerate}
	\item From a relative (10.6\%), 
	\item From a close friend (20.2\%),
	\item From an acquaintance (13.4\%)
	\item From a placement service (school, public or private) (13.1\%) 
	\item From an advertisement or a sign (16.4\%) 
	\item Directly contacted by employer (9.7\%) 
	\item Directly contacted the employer (in person or by phone) (16.6\%)
\end{enumerate}

Respondents coded as having answered ``Don't Know'' are set to missing on the outcome. About 44\% of respondents found out about their job from a social contact (either through relatives, close friends or acquaintances). However, and consistent with previous research and theory, among respondents who did use a social contact, respondents were more than twice as likely to rely on a strong tie (relatives or close friends) in comparison to using a weak tie.

\section{Results (General Social Survey)}
\subsection{Means-difference tests}
Are there differences in the extent to which persons engage in more extensive patterns of culture consumption for different respondent categories (divided according to the method through which they found out about their jobs)?  To answer this question, I regress the cultural variety scale against dummy-variable indicators of whether the respondent found out about their current job using the seven job-search methods shown in Table~\ref{tab1} (in addition to the ``don't know'' category).  The \emph{t}-test corresponding to whether the coefficient in this bivariate regression is significantly different from zero, is equivalent to a test of whether average number of cultural activities engaged in differ for a given group of respondents in comparison to the rest of the sample.  The results are shown in Table~\ref{tab1}.

What I find is broadly consistent with expectations.  Respondents who found out about their job through a strong tie (close friends) display narrower taste profiles than the average respondent ($t=-2.16, p<0.05$).  However, respondents who relied on contact information from a weak ties (e.g. an acquaintance) have broader than average cultural variety scores ($t=2.09, p<0.05$), as do respondents who were themselves contacted directly by the employer ($t=2.86, p<0.05$).   This reflects two different sources of network embeddedness, roughly consonant between the difference between ``outdegree'' (social capital as a result of knowing lots of others) and ``indegree'' (social capital as a result of being known by lots of others) in network analysis.  

In a similar way, respondents who used the non-relational method of calling or asking themselves or simply answering a sign or and ad display more restricted taste patterns than average, but neither of these differences reach conventional levels of statistical significance.  None of the other contrasts---except that which compares respondents who used an employment service to the rest---are statistically significant.   While the difference between respondents who learned about their current job through relatives (a version of strong ties) is in the expected direction, it is not within conventional levels of significance.  In addition, respondents who simply responded to an ad or sign (a non-relational method of search) should be less likely to consume a wide span of cultural forms, but the analysis shows that they are no different from the average respondent ($t=-0.35$). 

\subsection{Probit models with sample selection}
It is possible that the results shown in Table~\ref{tab1} are the spurious by-product of unobserved socio-demographic heterogeneity among persons with different cultural breadth profiles.  For instance, if the highly educated are more likely to use weak-ties to find a job (or to be directly contacted by the employer), and they are also more likely to consume culture then the difference in the cultural variety score between those who found out about their current job through and acquaintance and the rest of the sample might simply be due to differences in education and not the higher likelihood of those who consume a variety of culture to be better connected.  The same argument can be made via heterogeneity premised on age, gender or race.  

In addition, the statistics shown in the table suffer from an ``partial observability'' bias \citep{heckman78}:  information on the source of information regarding current job is only reliably observed from those currently employed ($N=740$). If unobserved individual characteristics that lead to employment are correlated with the unobserved individual characteristics that induce a person to have found out about their job through a particular pathway, then this might create spurious differences across different groups of respondents.  

I use probit models with sample selection \citep{heckman78} to account for both of these sources of spuriousness.  The selection equation consists of a probit regression predicting whether the individual is currently employed or not. This equation contains the usual predictors of employment: age, age-squared, highest degree attained, race, gender, marital status, urban residence and the presence of children in the household. The substantive equation then predicts the likelihood of having found out about the current job through one of the mechanisms listed in Table~\ref{tab1}.

Table~\ref{tab2} shows the results of a series of binary probit models with sample selection in which the dependent variable equals one if the respondent used one of the search methods considered above, and zero otherwise.  I use the cultural variety scale as a predictor, but estimate the effect while adjusting for the respondent's education, age, gender, status as a racial or ethnic minority (classification as either African-American or Hispanic) in the substantive equation.  I also control for the occupational prestige of the current occupation (relevant if members of high and low status occupations are disproportionately more likely to rely on certain job-search methods).

The results show that after adjusting for possible socio-demographic confounders and for possible sources of partial observability, the differential likelihood of those respondents who consume a wide variety of cultural goods to have obtained a job by reliance through particular ``socially embedded'' mechanisms remains.  The only result shown in Table~\ref{tab1} that appears to be due to heterogeneity separating high culture consumers from abstainers is that corresponding to using a strong tie to find a job.  After controlling for the relevant socio-demographic characteristics and for selection into the ranks of the employed, I find that high cultural variety respondents are no more or less likely to have found out about their job through a strong tie than persons with more restricted tastes. 

High-cultural variety respondents continue to be statistically more likely to have found about their job through a weak tie ($t=2.08$) or to have been directly contacted by their current employer ($t=2.34$). In addition, I find that adjusting for variation in socio-demographic characteristics brings to light a result that was {\em supressed} in Table~\ref{tab1} due to unobserved heterogeneity:  high cultural variety individuals are much less likely than chance to have found their job through the (non-relational) means of having called themselves ($t=-2.04$).

\section{Study 2:  2002 Survey for Public Participation in the Arts (CPS/SPPA)}
\subsection{2002 Survey for Public Participation in the Arts (SPPA)} 
The 2002 Survey of Public Participation in the Arts (SPPA) a periodic (fielded roughly every 4-6 years) survey of the non-institutionalized American population regarding arts and popular culture consumption practices and habits sponsored by the National Endowment for the Arts. For funding-related reasons 2002 the SPPA was administered by the U.S. Census Bureau as a supplement to the Current Population Survey (CPS). About 25\% of the households that were interviewed for the CPS also received questions from the SPPA.

\subsubsection{Cultural variety scale}
As a measure of cultural variety in tastes, I create a linear scale that is equal to the total number of cultural activities and practices that the respondent reports having engaged in the past 12 months. Respondents were queried about their musical tastes, their attendance of public performances of music, theater and dance, as well as their reading habits (reading non-fiction, plays, poetry and fiction) and whether they had attend the movies, a sports event or an amusement park, for a total of 38 items. The scale has good reliability (Cronbach's $\alpha=0.90$) showing that we are indeed tapping a single unidimensional construct (cultural breadth or variety).

The descriptive statistics (the average proportion of respondents who engaged in each activity) is shown on Table~\ref{tab2}. The table shows that the three most popular activities are, not surprising, going to the movies, reading a book and listening to rock music. The three most restricted activities are those associated with traditionally the fine performance arts and literature (Peterson and Kern 1996): reading a play, going to the opera and attending the ballet. For the subset of respondents that will form part of the empirical analysis---as detailed below---the mean number of cultural activities engaged in 9.3 with a standard deviation of 6.9.

\subsubsection{Network activation indicator} 
The 2002 SPPA survey was conducted as part the Current Population Survey. This means that detailed information is available on a variety of job, industry and occupation related outcomes for each respondent. The particular subset of respondents (N=540) that we are interested in for the purposes of this study are those that are currently unemployed but that are actively looking for a job.\footnote{ According to the U.S. Bureau of Labor Statistics, these individuals must meet the following four requirements: 1) they are adult civilians; 2) they did not hold a job during the preceding week; 3) they were available to take a job if one were offered; and 4) during the preceding month, they had sought work through specific "active" methods that put them in touch with employers.} As part of the CPS, those respondents who reported that the ``tried to find work within the last four weeks'' were asked the various means that they tried to find work.   Overall respondents identified 12 different ways in which they attempted to find a job in the last month. 

The distribution of respondent-choices (some respondents yield more than one observation since they report having attempted multiple strategies) across the job search strategy categories is shown in Table~\ref{tab3}. The table shows that, not surprisingly sending out a resume of filling out an application is the most prevalent way in which individual go about searching for employment, with close to 24\% of respondents reporting having used that method in the past month.  Two \emph{network based strategies} are shown in the table, (1) \emph{contacting the employer directly} (10.5\%), and (2) \emph{contacting friends and relatives} (10.9\%). For the purposes of constructing a measure of network activation,  I count the respondent as having \emph{activated a social tie} if he or she reports having used any of the two network-based strategies in the past month (about 21\% of respondents). Because persons are much more likely to have direct access to an employer through a weak tie \citep{yakubovich05}, I use having resorted to the \emph{direct contact strategy} as a \emph{indirect measure} of weak tie activation of the job-contact network. Finally I create a binary variable that equals one if the respondent reports having contacted family a friends in order to find a job, which I label a \emph{mediated contact strategy} as a measure of strong tie activation of the job-contact network.



\section{Results (Current Population Survey)}
\subsection{Means-difference tests}

\section{Discussion and Conclusion}
\subsection{Summary of the results}
The above results provide support for the proposition that the primary value of cultural variety in contemporary American society is \emph{social}: the ability to appropriate a wide variety of symbolic goods allows the individual to maintain---and diachronically replenish---\emph{a larger }and \emph{a more varied} personal network.   ``Omnivorous'' persons---when this is defined by the variety of tastes and cultural practices that persons report engaging in---are more likely to rely on socially embedded strategies when looking for a job.

This analysis contributes to our understanding of the role of culture in processes of stratification by connecting cultural capital not only to the static \emph{possession }of social resources in the form of larger and more varied networks of cultural relations, but by providing empirical evidence of a link between cultural capital and the \emph{agentic activation of those networks}. This is important because it is through the activation of resources embedded in social relations that social capital is theorized as going from a potential to a real resource \citep{lin99}. These results are consistent with the general notion that cultural taste variety ``counts'' as cultural capital in the contemporary societies precisely due to the intimate association of culture consumption with sociability, under conditions of ``blurring'' and ``deritualization'' of previously established boundaries across the arts.

\singlespacing
\newpage
\bibliographystyle{asr}
\bibliography{N:/Private/zotero-library}

%Figures
\newpage
\begin{figure}[htb!]
\caption{}
\centering
\includegraphics[width=6in, height=5in]{gss-graph.pdf}
\label{fig1}
\end{figure}


%Tables
\newpage
\begin{table}[htb!]
\caption{2002 General Social Survey}
\begin{tabular*}{\hsize}{@{\extracolsep{\fill}} l*{7}{c}} 
\hline
\hline
&\multicolumn{2}{c}{Mean N. Cultural Activities}&&\\
Found out about job through:&$\beta=0$&$\beta=1$&t-statistic&p-value\\
\hline
Relatives&2.49&2.40&-0.48&0.629\\
Close Friend&2.55&2.19&-2.46&0.014\\
Acquaintance&2.43&2.78&2.03&0.043\\
Emp. Service&2.46&2.67&1.14&0.255\\
Advertisement/Sign&2.51&2.33&-1.19&0.233\\
Directly From Employer&2.41&3.01&3.14&0.002\\
Called/Asked Myself&2.51&2.30&-1.29&0.196\\
\hline
\end{tabular*}
\label{tab1}
\end{table}

\newpage
\begin{landscape}
\begin{table}[htb!]
\caption{\footnotesize 2002 General Social Survey.  (1) Relatives , (2) Close Firend, (3) Acquaintance, (4) Emp. Service/School , (5) Advertisement/Sign , (6) Direct Contact with Employer, (7) Called/Asked Myself.}
\begin{tabular*}{\hsize}{@{\extracolsep{\fill}} l*{7}{c}} 
\hline
\hline 
&(1)&(2)&(3)&(4)&(5)&(6)&(7)\\
\hline
Cultural Variety Scale   &     0.06&    -0.06&     0.16&    -0.05&    -0.12&     0.22&    -0.15\\
                         &   (0.68)&  (-0.85)&   (2.19)&  (-0.54)&  (-1.62)&   (2.66)&  (-1.95)\\
N. of Years of Education &    -0.02&    -0.04&    -0.04&     0.04&    -0.03&     0.05&     0.07\\
                         &  (-0.44)&  (-0.89)&  (-0.79)&   (0.71)&  (-0.57)&   (0.84)&   (1.45)\\
Occ. Prestige (Current Job)&    -0.03&    -0.01&    -0.00&     0.03&     0.02&    -0.01&     0.00\\
                         &  (-2.99)&  (-1.75)&  (-0.37)&   (3.27)&   (1.90)&  (-0.62)&   (0.27)\\
Respondent's Age         &    -0.04&     0.00&    -0.00&    -0.01&    -0.01&     0.04&     0.01\\
                         &  (-3.57)&   (0.00)&  (-0.43)&  (-0.96)&  (-0.95)&   (4.10)&   (1.66)\\
Gender (Women=1)         &    -0.26&    -0.20&     0.10&    -0.02&     0.11&     0.23&     0.08\\
                         &  (-1.02)&  (-1.02)&   (0.45)&  (-0.07)&   (0.53)&   (0.90)&   (0.38)\\
Race/Eth. (Black or Hispanic=1)&    -1.02&     0.56&    -0.11&     0.21&    -0.31&     0.42&    -0.03\\
                         &  (-2.35)&   (2.13)&  (-0.36)&   (0.62)&  (-1.04)&   (1.16)&  (-0.11)\\
Mar. Status (Married=1)  &     0.04&     0.63&     0.19&    -0.25&    -0.37&     0.39&    -0.61\\
                         &   (0.16)&   (3.01)&   (0.80)&  (-0.94)&  (-1.69)&   (1.46)&  (-2.63)\\
N. of Children           &     0.02&    -0.14&     0.01&    -0.07&     0.18&    -0.10&     0.06\\
                         &   (0.17)&  (-1.75)&   (0.14)&  (-0.67)&   (2.31)&  (-1.01)&   (0.76)\\
City Size in 1000s (logged)&     0.02&    -0.02&     0.07&    -0.03&     0.06&    -0.01&    -0.11\\
                         &   (0.35)&  (-0.37)&   (1.19)&  (-0.48)&   (1.18)&  (-0.08)&  (-1.87)\\ \\
Model $\chi^{2}$                       &    38.64&    27.04&     7.54&    22.51&    14.00&    31.71&    18.51\\
N                        &      717&      717&      717&      717&      717&      717&      717\\
\hline
\end{tabular*}
\label{tab2}
\end{table}


\end{landscape}

\end{document}

















